\documentclass[10pt]{beamer}

\usetheme[bib=biblatex]{enlp}
\enlpusebibresource{references.bib}

\enlpsetpagesize{33.87cm}{19.05cm}
\enlpsetsubtitleposition{.063\paperwidth}{-.097\paperheight}
\enlpsetdateposition{.063\paperwidth}{-.145\paperheight}
\enlpsettitleposition{.50\paperwidth}{-.43\paperheight}
\enlpsetinstituteposition{.50\paperwidth}{-.64\paperheight}
\enlpsetauthorposition{.50\paperwidth}{-.72\paperheight}
\enlpsetsubtitlewidth{.36\paperwidth}
\enlpsetdatewidth{.30\paperwidth}
\enlpsettitlewidth{.78\paperwidth}
\enlpsetinstitutewidth{.78\paperwidth}
\enlpsetauthorwidth{.78\paperwidth}

\title{ENLP Beamer Theme Documentation}
\subtitle{User Guide and Package Reference}
\author{Eito Yoneyama}
\institute{Ehime University Natural Language Processing Laboratory}
\date{\today}

\begin{document}

\enlptitleframe

\enlpsectionframe{Overview}

\begin{frame}{Purpose}
  \begin{itemize}
    \item This document is the primary user documentation for the ENLP Beamer theme.
    \item The package reproduces the ENLP PowerPoint template as a Beamer theme with natively drawn Ti\textit{k}Z backgrounds.
    \item The official theme name is \texttt{\string\usetheme\{enlp\}}.
    \item The document is also built by \texttt{l3build}, so the rendered PDF doubles as a packaging check.
  \end{itemize}
\end{frame}

\begin{frame}{Requirements}
  \begin{itemize}
    \item LuaLaTeX, XeLaTeX, pLaTeX, and upLaTeX are supported. pdfLaTeX is not supported.
    \item pLaTeX and upLaTeX must use \texttt{dvipdfmx}, i.e.\ \texttt{\string\documentclass[dvipdfmx]\{beamer\}}.
    \item The default page size is 33.87cm $\times$ 19.05cm.
    \item The default Latin sans font is \texttt{TeX Gyre Heros}; the default Japanese sans font is \texttt{HaranoAji Gothic}.
    \item The default monospaced font is \texttt{Latin Modern Mono}.
    \item Under LuaLaTeX, mixed Japanese-Latin text is handled by \texttt{luatexja-fontspec} on top of \texttt{luaotfload}.
    \item Title, heading, and content backgrounds are drawn natively using Ti\textit{k}Z.
  \end{itemize}
\end{frame}

\begin{frame}{Engine-Specific Notes}
  \begin{itemize}
    \item LuaLaTeX uses \texttt{fontspec} with \texttt{luatexja-fontspec}; \texttt{luaotfload} resolves the Latin and Japanese OpenType fonts.
    \item XeLaTeX uses \texttt{fontspec} with \texttt{xeCJK} and the same HaranoAji font family.
    \item pLaTeX and upLaTeX use a \texttt{dvipdfmx}-based workflow with \texttt{pxchfon} and \texttt{pxjahyper}.
    \item The theme defaults both Latin and Japanese body text to sans serif so mixed Japanese-Latin slides stay visually consistent.
    \item Hyperlinks and bookmarks are configured for Unicode output across every supported engine.
  \end{itemize}
\end{frame}

\begin{frame}{Package Contents}
  \begin{itemize}
    \item \texttt{beamerthemeenlp.sty}: main theme integration
    \item \texttt{beamer*themeenlp.sty}: split sub-themes (color, font, inner, outer)
    \item \texttt{enlplogo.sty}: lab logo definitions
    \item \texttt{enlpbeamer.bbx}, \texttt{enlpbeamer.cbx}, \texttt{enlpbeamer.bst}: bibliography styles
  \end{itemize}
\end{frame}

\enlpsectionframe{Getting Started}

\begin{frame}{Minimal Preamble for LuaLaTeX / XeLaTeX}
\begin{block}{Basic Usage}
\texttt{\string\documentclass\{beamer\}}\\
\texttt{\string\usetheme\{enlp\}}\\
\texttt{\string\title\{Title\}}\\
\texttt{\string\author\{Author\}}\\
\texttt{\string\institute\{Institute\}}\\
\texttt{\string\date\{\string\today\}}
\end{block}
  \begin{itemize}
    \item The theme applies the ENLP page size automatically.
    \item No manual background configuration is required.
    \item Japanese and Latin text are both initialized to sans-serif defaults for stable mixed-script slides.
    \item The recommended title slide command is \texttt{\string\enlptitleframe}.
  \end{itemize}
\end{frame}

\begin{frame}{Minimal Preamble for pLaTeX / upLaTeX}
\begin{block}{Basic Usage}
\texttt{\string\documentclass[dvipdfmx]\{beamer\}}\\
\texttt{\string\usetheme\{enlp\}}\\
\texttt{\string\title\{Title\}}\\
\texttt{\string\author\{Author\}}\\
\texttt{\string\institute\{Institute\}}\\
\texttt{\string\date\{\string\today\}}
\end{block}
  \begin{itemize}
    \item The document class must use the \texttt{dvipdfmx} option.
    \item The same mixed Japanese-Latin sans-serif defaults are applied after the theme loads.
    \item The recommended title slide command is \texttt{\string\enlptitleframe}.
  \end{itemize}
\end{frame}

\begin{frame}{Recommended Document Skeleton}
\begin{block}{Typical Structure}
\texttt{\string\begin\{document\}}\\
\texttt{\string\enlptitleframe}\\
\texttt{\string\enlpsectionframe\{Introduction\}}\\
\texttt{\string\begin\{frame\}\{Motivation\}}\\
\texttt{\ \ ...}\\
\texttt{\string\end\{frame\}}\\
\texttt{\string\insertbibliographyframe}\\
\texttt{\string\end\{document\}}
\end{block}
\end{frame}

\begin{frame}{Custom Fonts Overview}
  \begin{itemize}
    \item Custom fonts can already be configured in the preamble after \texttt{\string\usetheme\{enlp\}}.
    \item LuaLaTeX uses the \texttt{fontspec} and \texttt{luatexja-fontspec} interfaces.
    \item XeLaTeX uses the \texttt{fontspec} and \texttt{xeCJK} interfaces.
    \item pLaTeX and upLaTeX use \texttt{pxchfon} for Japanese fonts and regular NFSS family selection for Latin fonts.
    \item Because the theme defaults to sans-serif and gothic text, update the sans and Japanese gothic side when changing body fonts.
  \end{itemize}
\end{frame}

\begin{frame}{Custom Fonts with LuaLaTeX}
\begin{block}{Example}
\small
\texttt{\string\usetheme\{enlp\}}\\
\texttt{\string\setmainfont\{TeX Gyre Termes\}}\\
\texttt{\string\setsansfont\{IBM Plex Sans\}}\\
\texttt{\string\setmonofont\{JetBrains Mono\}}\\
\texttt{\string\setmainjfont[BoldFont=HaranoAjiMincho-Bold.otf]\{HaranoAjiMincho-Regular.otf\}}\\
\texttt{\string\setsansjfont[BoldFont=HaranoAjiGothic-Bold.otf]\{HaranoAjiGothic-Regular.otf\}}
\end{block}
  \begin{itemize}
    \item Place these commands after \texttt{\string\usetheme\{enlp\}} and before \texttt{\string\begin\{document\}}.
    \item Use \texttt{\string\setmainjfont} and \texttt{\string\setsansjfont} for the Japanese side of mixed Japanese-Latin text.
    \item The theme uses sans-serif and Japanese gothic text for slide bodies, so \texttt{\string\setsansfont} and \texttt{\string\setsansjfont} are the pair that changes what you see on most frames.
    \item LuaLaTeX relies on \texttt{luatexja-fontspec} for Japanese layout and \texttt{luaotfload} for font resolution, so this is the preferred path when mixed Japanese-Latin spacing matters.
    \item If you need luatexja-specific options such as JFM tuning, pass them in the optional argument.
  \end{itemize}
\end{frame}

\begin{frame}{LuaLaTeX Font Troubleshooting}
\begin{block}{If \texttt{\string\setmainfont} seems ignored}
\small
\texttt{\string\usetheme\{enlp\}}\\
\texttt{\string\setsansfont[BoldFont=*Bold,FakeBold=2.5]\{LINE Seed JP\}}\\
\texttt{\string\setsansjfont[BoldFont=*Bold,FakeBold=2.5]\{LINE Seed JP\}}
\end{block}
  \begin{itemize}
    \item This theme is not like \texttt{ltjsarticle}: normal slide text does not use the Roman family by default.
    \item \texttt{\string\setmainfont} and \texttt{\string\setmainjfont} only affect the Roman side, so they may not visibly change the frame body.
    \item If you write both \texttt{\string\setmainfont[\ldots]\{...\}} and later \texttt{\string\setmainfont\{...\}}, the later command replaces the earlier options.
    \item Loading \texttt{luatexja} and \texttt{luatexja-fontspec} manually is optional because the theme already loads them under LuaLaTeX.
  \end{itemize}
\end{frame}

\begin{frame}{Custom Fonts with XeLaTeX}
\begin{block}{Example}
\small
\texttt{\string\usetheme\{enlp\}}\\
\texttt{\string\setmainfont\{TeX Gyre Termes\}}\\
\texttt{\string\setsansfont\{IBM Plex Sans\}}\\
\texttt{\string\setmonofont\{JetBrains Mono\}}\\
\texttt{\string\setCJKmainfont[BoldFont=HaranoAjiMincho-Bold.otf]\{HaranoAjiMincho-Regular.otf\}}\\
\texttt{\string\setCJKsansfont[BoldFont=HaranoAjiGothic-Bold.otf]\{HaranoAjiGothic-Regular.otf\}}
\end{block}
  \begin{itemize}
    \item XeLaTeX uses \texttt{xeCJK}, so Japanese fonts are changed with \texttt{\string\setCJKmainfont} and \texttt{\string\setCJKsansfont}.
    \item The theme keeps body text in sans-serif by default, so \texttt{\string\setsansfont} and \texttt{\string\setCJKsansfont} are the most important pair.
    \item Font names and font file names are both acceptable if XeTeX can resolve them.
  \end{itemize}
\end{frame}

\begin{frame}{Custom Fonts with pLaTeX / upLaTeX}
\begin{block}{Example}
\small
\texttt{\string\documentclass[dvipdfmx]\{beamer\}}\\
\texttt{\string\usepackage\{tgtermes\}}\\
\texttt{\string\usetheme\{enlp\}}\\
\texttt{\string\renewcommand\{\string\rmdefault\}\{qtm\}}\\
\texttt{\string\renewcommand\{\string\sfdefault\}\{qhv\}}\\
\texttt{\string\renewcommand\{\string\ttdefault\}\{lmtt\}}\\
\texttt{\string\setminchofont\{HaranoAjiMincho-Regular.otf\}}\\
\texttt{\string\setboldminchofont\{HaranoAjiMincho-Bold.otf\}}\\
\texttt{\string\setgothicfont\{HaranoAjiGothic-Regular.otf\}}\\
\texttt{\string\setboldgothicfont\{HaranoAjiGothic-Bold.otf\}}
\end{block}
  \begin{itemize}
    \item Latin fonts are changed through normal LaTeX family declarations such as \texttt{\string\rmdefault}, \texttt{\string\sfdefault}, and \texttt{\string\ttdefault}.
    \item Japanese fonts are changed through \texttt{pxchfon} commands such as \texttt{\string\setminchofont} and \texttt{\string\setgothicfont}.
    \item Keep \texttt{dvipdfmx} enabled or the theme backgrounds and Japanese font mapping will not work correctly.
    \item The same procedure applies to both pLaTeX and upLaTeX; upLaTeX is usually preferable for fully Unicode sources.
  \end{itemize}
\end{frame}

\begin{frame}{Default Typography}
  \begin{columns}[T,totalwidth=\textwidth]
    \begin{column}{0.48\textwidth}
      \begin{block}{Title and Heading Slides}
        \begin{itemize}
          \item title: 44pt
          \item subtitle: 20pt
          \item author / institute / date: 24pt
          \item section title: 48pt
        \end{itemize}
      \end{block}
    \end{column}
    \begin{column}{0.48\textwidth}
      \begin{block}{Regular Slides}
        \begin{itemize}
          \item frame title: 36pt
          \item body text: 24pt
          \item block title / block body: 24pt
          \item footnote / footline: 12pt
        \end{itemize}
      \end{block}
    \end{column}
  \end{columns}
\end{frame}

\enlpsectionframe{Layout Control}

\begin{frame}{Title Slide Defaults}
  \begin{itemize}
    \item The title slide background is drawn internally via the outer theme.
    \item The subtitle and date are placed in the upper-left metadata region.
    \item The title is centered in the large middle box.
    \item Institute and author are centered in the lower metadata region.
  \end{itemize}
\end{frame}

\begin{frame}{Title Slide Position Macros}
\begin{block}{Position Controls}
\texttt{\string\enlpsetsubtitleposition\{<x>\}\{<y>\}}\\
\texttt{\string\enlpsetdateposition\{<x>\}\{<y>\}}\\
\texttt{\string\enlpsettitleposition\{<x>\}\{<y>\}}\\
\texttt{\string\enlpsetinstituteposition\{<x>\}\{<y>\}}\\
\texttt{\string\enlpsetauthorposition\{<x>\}\{<y>\}}
\end{block}
\begin{block}{Width Controls}
\texttt{\string\enlpsetsubtitlewidth\{<width>\}}\\
\texttt{\string\enlpsetdatewidth\{<width>\}}\\
\texttt{\string\enlpsettitlewidth\{<width>\}}\\
\texttt{\string\enlpsetinstitutewidth\{<width>\}}\\
\texttt{\string\enlpsetauthorwidth\{<width>\}}
\end{block}
\end{frame}

\begin{frame}{Title Slide Default Coordinates}
  \begin{itemize}
    \item subtitle: \texttt{x=.063w}, \texttt{y=-.090h}, \texttt{width=.28w}
    \item date: \texttt{x=.063w}, \texttt{y=-.145h}, \texttt{width=.30w}
    \item title: \texttt{x=.50w}, \texttt{y=-.47h}, \texttt{width=.78w}
    \item institute: \texttt{x=.50w}, \texttt{y=-.64h}, \texttt{width=.78w}
    \item author: \texttt{x=.50w}, \texttt{y=-.72h}, \texttt{width=.78w}
  \end{itemize}
\end{frame}

\begin{frame}{Regular Slide Geometry}
  \begin{itemize}
    \item Regular slides use the natively drawn content background.
    \item The default left and right text margins are both 24mm.
    \item The content width is matched to the title underline length in the background.
    \item Frame contents are top-aligned instead of vertically centered.
    \item The frame number is placed in the lower right; custom footer text appears in the lower left.
  \end{itemize}
\end{frame}

\begin{frame}{Block and List Styling}
  \begin{itemize}
    \item Blocks use an ENLP-specific template sized for 24pt body text.
    \item Block titles are white on a dark turquoise background.
    \item Block bodies use a pale wash background.
    \item Item markers are enlarged slightly relative to Beamer defaults for better visibility at 24pt.
  \end{itemize}
\end{frame}

\enlpsectionframe{Section Slides}

\begin{frame}{Section Heading Pages}
  \begin{itemize}
    \item Section heading pages use the natively drawn heading background.
    \item The recommended command is \texttt{\string\enlpsectionframe\{...\}}.
    \item This command registers the section and immediately inserts the heading page.
    \item The default heading size is 48pt.
  \end{itemize}
\end{frame}

\begin{frame}{Optional Section Subtitle}
\begin{block}{Usage}
\texttt{\string\sectionpagesubtitle\{Spring Semester\}}\\
\texttt{\string\enlpsectionframe\{Background\}}
\end{block}
  \begin{itemize}
    \item The subtitle is shown only on the section heading page.
    \item Pass an empty string to clear it for later sections.
  \end{itemize}
\end{frame}

\enlpsectionframe{Bibliography}

\begin{frame}{Supported Backends}
  \begin{itemize}
    \item The theme supports both \texttt{biblatex + biber} and \texttt{BibTeX}.
    \item The theme option \texttt{bib=biblatex} is the default.
    \item \texttt{bib=bibtex} switches to the bundled \texttt{enlpbeamer.bst}.
    \item In either mode, bibliography files are registered with \texttt{\string\enlpusebibresource\{...\}}.
    \item If a document does not use bibliography features yet, \texttt{\string\enlpusebibresource} may be omitted entirely.
  \end{itemize}
\end{frame}

\begin{frame}{Backend Selection}
\begin{block}{biblatex + biber}
\texttt{\string\usetheme[bib=biblatex]\{enlp\}}\\
\texttt{\string\enlpusebibresource\{references.bib\}}
\end{block}
\begin{block}{BibTeX}
\texttt{\string\usetheme[bib=bibtex]\{enlp\}}\\
\texttt{\string\enlpusebibresource\{references.bib\}}
\end{block}
\begin{block}{Backend Initialization Without a .bib File}
\texttt{\string\usetheme[bib=biblatex]\{enlp\}}\\
\texttt{\string\enlpusebibresource\{\}}
\end{block}
\begin{itemize}
  \item This is useful when you want to decide the backend up front but do not have a bibliography database yet.
  \item The same empty registration also works with \texttt{bib=bibtex}.
\end{itemize}
\end{frame}

\begin{frame}{Citation Semantics}
  \begin{itemize}
    \item \texttt{\string\cite\{...\}} and \texttt{\string\enlptextcite\{...\}} produce superscripted numeric references such as \cite{doe2025}.
    \item \texttt{\string\enlpauthorcite\{...\}} produces an author-style citation, for example \enlpauthorcite{smith2024}.
    \item \texttt{\string\enlpparencite\{...\}} produces a parenthetical citation.
    \item Citation keys used inside a frame are collected automatically and rendered as full entries in the footer.
    \item Footer bibliography labels remain ordinary bracketed numbers such as \texttt{[1]}.
  \end{itemize}
\end{frame}

\begin{frame}{Automatic Footer Citations}
  This frame cites \cite{doe2025} and \enlpparencite{smith2024}, so both entries appear in the footer automatically.
  \begin{itemize}
    \item Repeated citations are deduplicated in first-use order.
    \item \texttt{\string\usetheme[autocitefooter=false]\{enlp\}} disables the automatic footer behavior.
  \end{itemize}
\end{frame}

\begin{frame}{Manual Footer Citations and Bibliography Pages}
  \framefullcite{doe2025,smith2024}
  \begin{itemize}
    \item \texttt{\string\framefullcite\{...\}} overrides the automatic footer content with a fixed citation set.
    \item Footer entries are rendered in bibliography-number order, even if keys are passed in a different order.
    \item Footer entries use the 12pt footer font, not the 24pt body font.
    \item \texttt{\string\insertbibliographyframe} inserts a bibliography section heading followed by bibliography pages.
    \item Bibliography pages themselves use the same 24pt body layout as regular content slides.
  \end{itemize}
\end{frame}

\begin{frame}{Bibliography Color}
  \begin{itemize}
    \item The default bibliography color is black.
    \item This applies to in-text citations, footer citations, and bibliography pages.
    \item Use \texttt{\string\enlpsetbibliographycolor\{...\}} to override it globally.
    \item \texttt{\string\enlpsetcitecolor\{...\}} is provided as an alias.
  \end{itemize}
\end{frame}

\enlpsectionframe{Build and Install}

\begin{frame}{l3build Workflow}
  \begin{itemize}
    \item \texttt{l3build doc} builds this documentation.
    \item \texttt{l3build install --full --dry-run} verifies the TDS installation layout without copying files into the live tree.
    \item When running LuaLaTeX in a sandboxed or CI environment, make sure \texttt{luaotfload} has writable \texttt{TEXMFVAR} and \texttt{TEXMFCACHE} paths.
    \item The documentation PDF and source are installed under \texttt{doc/latex/enlp-beamer/}.
    \item The theme files, logo package, and bibliography styles are installed under \texttt{tex/latex/enlp-beamer/}.
  \end{itemize}
\end{frame}

\begin{frame}{TDS Layout}
  \begin{itemize}
    \item \texttt{tex/latex/enlp-beamer/beamerthemeenlp.sty}
    \item \texttt{tex/latex/enlp-beamer/beamer*themeenlp.sty}
    \item \texttt{tex/latex/enlp-beamer/enlplogo.sty}
    \item \texttt{bibtex/bst/enlp-beamer/enlpbeamer.bst}
    \item \texttt{doc/latex/enlp-beamer/documentation.pdf}
  \end{itemize}
\end{frame}

\begin{frame}{Design Policy}
  \begin{itemize}
    \item Background assets are defined directly using Ti\textit{k}Z within the outer theme.
    \item User-facing customization is focused on text placement, page size, bibliography backend, and bibliography color.
    \item This keeps visual reproduction stable across local builds and installed TDS trees.
    \item If the underlying artwork changes, the Ti\textit{k}Z drawing commands in the outer theme should be updated.
  \end{itemize}
\end{frame}

\enlpsectionframe{References}
\insertbibliographyframe*[References]

\end{document}
